\subsection{Large corners with four objects and sixteen arrows}

\begin{lemma}\label{C:three-small-H2-zero}
A connected three-object category with at most seven arrows has
$H^2(BD;\mathbb Q)=0$. The same holds without connectedness at
at most six arrows.
\end{lemma}
\begin{proof}
In the connected case a corner has order at most three, since
there are two other identities and at least two cross arrows.
A retract deletion leaves at most four arrows, whose rational
$H^2$ is zero by the complete small-category census. Assume
no retract is present. A nontrivial strong pair without
such a retract requires two corners of order at least two,
two opposite cross arrows, and one further connecting cross
arrow, hence at least eight arrows. Thus the support is acyclic.
If a corner has order three, the remaining two cross arrows
are singletons and the other corners are trivial. There is
no composable pair of cross arrows, since that would require
a third cross hom-set. The nerve is a graph of vertex spaces
and singleton correspondence spaces. All vertex monoids have
rationally acyclic classifying spaces, by the complete order-four
classification; a singleton correspondence is their product.
The graph is a tree, so Mayer--Vietoris gives rational acyclicity.

Otherwise all corner algebras are rationally semisimple. If the
support has no three-object chain, its relative bar has length
at most one. If it has such a chain, write its three hom-set
sizes as $x,y,z$, with $z$ the long edge. If some corner has
order two, all three hom-sets are singletons, and one of the
trivial extreme corners is initial or terminal. If all corners
are trivial, $x+y+z\leq4$. The possibilities $x=2,y=z=1$
and $y=2,x=z=1$ have a terminal or initial object. In the other
possibilities $x=y=1$, the single two-simplex has a nonzero
boundary, so $H^2=0$.

For a disconnected three-object category with at most six
arrows, an isolated component monoid has order at most four.
A connected two-object component has at most five arrows. A nontrivial strong component either has a trivial
corner, which is a retract, or needs at least six arrows.
A triangular component with nontrivial corner and two cross
arrows has corner orders $(2,1)$; the order-two coinvariant
functor is exact and its orbit-collapse gives a graph.
All remaining triangular components have singleton cross hom
or both corners trivial. A singleton gives the join of the
vertex spaces, which is rationally acyclic here, and trivial
corners give a graph. These cases prove the assertion.
\end{proof}

\begin{proposition}\label{C:sixteen-four-large-corner}
A four-object category with at most sixteen arrows and an
endomorphism monoid of order at least eight cannot have type
$\mathbb{CP}^2$.
\end{proposition}
\begin{proof}
Assume the category is connected, and denote the large corner
by $v$, its order by $q\geq8$. A retract deletion leaves at
most thirteen arrows, already excluded. We may assume there
is no retract. Connectivity bounds $q\leq10$.

First exclude a nontrivial strong component. Every object in
one has a nontrivial corner, since otherwise it is a retract
of another object. The only possible budget is a strong pair
containing $v$, with corner orders $8,2$, two other trivial
corners, and exactly four cross arrows. Every nonempty cross
hom-set is a singleton. At the small corner the composite of
the opposite cross arrows is its nonidentity zero idempotent.
It fixes every incoming arrow. Lemma~\ref{C:one-sided-corner-deletion}
deletes that object and again reduces to at most thirteen arrows.
Thus the support can be assumed acyclic.

Suppose $v$ has both incoming and outgoing cross arrows. A
three-object chain through $v$ and its composite, together
with a connection to the fourth object, require at least
four cross arrows. If there are exactly four, the total
incoming and outgoing counts at $v$ must each be one:
a sum at least three requires at least five cross arrows,
whether the three incident arrows meet two or three neighbors.
The two adjacent corners are distinct, and at most one can
have order two; the other is trivial. Thus one of the two
singleton cofinal deletions deletes $v$.
If there are five cross arrows, $q=8$ and every other corner
is trivial. Both the incoming and outgoing total cannot be
at least two. With only two neighbors this would require
four incident arrows, a composite, and a further connecting
arrow; with three neighbors it would require four incident
arrows and at least two different composite hom-sets.
Either way it requires at least six arrows. Hence one of the
incoming and outgoing totals is one, and the adjacent corner
is trivial. The same cofinal deletion applies. Its result
has at most seven arrows. We may therefore assume that $v$
is extreme, and use opposition to make it a sink.

Let $D=C\setminus\{v\}$. If any incoming hom-set at $v$ is a
singleton, then $D$ has at most seven arrows. If it has seven,
the budget forces $q=8$ and exactly one cross arrow incident
to $v$, so connectivity forces $D$ to be connected. If it has
at most six, connectedness is unnecessary.
Lemma~\ref{C:three-small-H2-zero} gives $H^2(BD;\mathbb Q)=0$.
The singleton incoming arrow is fixed by every element of the
large monoid. Lemma~\ref{C:fixed-incoming-cups} therefore kills
all rational degree-two products.

It remains to suppose every nonempty incoming hom-set has at
least two elements. If there are at least two neighbors of
$v$, there are exactly two, each with two arrows to $v$;
$q=8$, all other corners are trivial, and one singleton cross
arrow connects the fourth object $w$ to one of these neighbors
$x$. It must be $x\to w$: the other orientation would force
a third incoming hom-set to $v$. Its sole incoming arrow
comes from the trivial corner $x$, so the singleton cofinal
deletion removes $w$, leaving at most fourteen arrows.

Finally suppose there is just one neighbor $w$, with
$s=|\operatorname{Hom}(w,v)|\geq2$. There are no cross arrows
into $w$ from $D\setminus\{w\}$, because their composites
would give another neighbor of $v$. Thus the nerve is the
union of $ND$ and the full two-object nerve $NY$ on $w,v$,
intersecting in $N\End(w)$; this is a homotopy pushout of
subcomplexes. The category $D$ is connected. Its two necessary
cross arrows give $s\leq3$ and $|\End(w)|\leq2$.
For $s=2$, $D$ has at most six arrows, so its rational $H^2$
is zero. The common corner is rationally acyclic. The
Mayer--Vietoris sequence consequently identifies $H^2(BC;\mathbb Q)$
with $H^2(BY;\mathbb Q)$ and injects $H^4(BC;\mathbb Q)$
into $H^4(BD;\mathbb Q)\oplus H^4(BY;\mathbb Q)$.
Lemma~\ref{C:two-object-cup-pruning} kills the degree-two
products on $BY$. Their restrictions on $BD$ are also zero,
so their products on $BC$ vanish.
For $s=3$, the entire budget forces $q=8$, every corner of
$D$ trivial, and exactly two singleton cross arrows in $D$.
Its support is a tree, so $BD$ is contractible, and the
pushout gives $BC\simeq BY$. This is a three-state act cone
with twelve arrows, excluded by
Theorem~\ref{C:three-state-cone-minimum22}.
Every case contradicts the primitive nonzero projective-plane
square or an already proved smaller-category exclusion.
\end{proof}
